\documentclass[../main]{subfiles}
\begin{document}

\chapter{Máquina Rotativa (MCU)}
\img{images/10/motor1.png}{Máquina Rotativa}
\info{
Movimiento Circular Uniforme}{El movimiento circular uniforme (M.C.U.) es el que describe un cuerpo que se mueve alrededor de un eje de giro con un radio y una velocidad angular $\omega$ constantes, trazando una circunferencia y con una aceleración centrípeta $a_c$. En este movimiento la dirección varia en cada instante, un ejemplo de este movimiento es un motor eléctrico o un alternador que gira a una $\omega$ constante.}

\info{Frecuencia mecánica}{
La frecuencia (f), es el número de revoluciones en la unidad de tiempo. 
\begin{equation}
f=\dfrac{vueltas}{t}
\end{equation}

siendo:

$f:$Frecuencia $[1/s][Hz] (Hertz)$

$t:$ Tiempo en segundos $[s]$
}
\info{Período}{
El Período $T$, de un cuerpo en movimiento circular uniforme es el tiempo empleado en efectuar una vuelta completa o revolución.
\begin{equation}
 T=\dfrac{1}{f}   
\end{equation}

siendo:

$f$: Frecuencia $[Hz]$

$T$: Período $[s]$
}
\info{RPM}{ \textbf{R}evoluciones \textbf{P}or \textbf{M}inuto, son las vueltas que se realizan en un minuto, su unidad es RPM.
\begin{equation}
 n= 60 f   
\end{equation}

siendo:

$n:$ revoluciones por minuto $[RPM]$

$f:$ frecuencia.

}
\info{Velocidad o rapidez Angular}{
\emph{Rapidez angular}

$\omega=2\pi f=\dfrac{2\pi}{T}=\dfrac{\Delta \vartheta}{\Delta t} =\dfrac{\vartheta_f-\vartheta_i}{ t_f - t_i}=\dfrac{2 \times \pi \times n}{60}$

}







\actividad{Resolver los siguientes problemas.}

\begin{enumerate}
    \item Un motor se mueve con velocidad $n=3000 RPM$, calcular su frecuencia $f$ en Hertz.
    \item Calcular el período $T$ y la velocidad angular $\omega$ del problema anterior.
    \item Un motor gira con movimiento circular uniforme, girando con $\omega=7\pi \dfrac{rad}{s}$. Calcular el ángulo que barre el radio de giro en $\Delta t=2s$. 
        \item Una bobina realiza un movimiento circular uniforme, da 3 vueltas en 1 minuto alrededor de su centro de giro. Calcular su rapidez angular (en rad/s).

    \item El motor de la imagen realiza un M.C.U. con
una rapidez angular de $\omega=2\pi \dfrac{rad}{s}$ .Determine el
tiempo que emplea para ir desde A hasta B. 
    \begin{center}
    \begin{tikzpicture}
    \draw [dashed](0,0)circle (2);
    \draw (0,0)--(-30:2) node [right]{$B$};
    \draw (0,0)--(90:2) node [above]{$A$};
    \draw (1,0.5) node {$\dfrac{2\pi}{3}$};
    \end{tikzpicture}
    \end{center}
\item Los cables A y B giran con MCU, con velocidad angulares de $\omega_A=3rad/s$ y $\omega_B=5rad/s$ respectivamente. Calcular el tiempo que tarde el cable B en alcanzar al cable A, tomando en cuenta el gráfico. 
    \begin{center}
    \begin{tikzpicture}
    \draw [dashed](0,0)circle (2);
    \draw (0,0)--(-90:2) node [below]{$B$};
    \draw (0,0)--(90:2) node [above]{$A$};
    \draw (1,0) node {$\pi$};
    \end{tikzpicture}
    \end{center}
\info{Velocidad o rapidez tangencial}{

\begin{equation}
v_t=\omega \times r
\end{equation}

Siendo:

$v_t$: Velocidad tangencial $[m/s]$.

$\omega:$ Velocidad angular $[rad/s]$

$r:$ Radio. [m]
}


    
    \info{Aceleración centrípeta}{
\begin{equation}
a_c=\omega^2 \times r = \dfrac{v_t^2}{r}
\end{equation}

Siendo:

$a_c$: Aceleración centrípeta $[m/s^2]$.

$v_t$: Velocidad tangencial $[m/s]$.

$\omega:$ Velocidad angular $[rad/s]$

$r:$ Radio. [m]
}


    
    \item La rapidez tangencial de una bobina que gira con MCU es de $v_t=12 m/s$. Calcular su aceleración centrípeta, si su radio mide $r=120 cm$.

    \item La frecuencia de una bobina que gira con M.C.U. es de $f=0,5Hz$. Hallar la rapidez tangencial en la periferia del bobinado, si tiene un diámetro de $d=40cm$.
    

    
\item Los puntos periféricos de un rotor jaula de ardilla con M.C.U se mueven a razón de $v_t=0,8 m/s$. Los puntos que están a 4cm de la periferia se
mueven a razón de $v_t=0,6 m/s$. Hallar el radio rotor. 

\item Un motor gira con MCU. Si los puntos
periféricos tienen el triple de rapidez tangencial
que aquellos puntos que se encuentran 5 cm
más cerca del centro motor, calcular el radio
del rotor. 

\end{enumerate}
\newpage
\end{document}
